% ==========================================================================
%  NOTE: This is NOT the actual thesis. Working notes / scratch material on an
%  EXPERIMENT: "how well can an AI agent integrate an existing scientific Shiny
%  app (BioDT Honeybee) into the ShinySwarm architecture?". Same register as
%  biodt_feasibility.tex -- raw notes to review and selectively pull into the
%  real thesis later (Methods / Results / Discussion). Verify every claim
%  against the live repos and the committed code before citing. Session run on
%  2026-06-12 with an AI coding agent (Ona) inside the project devcontainer.
% ==========================================================================

\documentclass[12pt,a4paper]{article}
\usepackage[utf8]{inputenc}
\usepackage[T1]{fontenc}
\usepackage{geometry}
\geometry{margin=2.5cm}
\usepackage{hyperref}
\usepackage{booktabs}
\usepackage{enumitem}
\usepackage{xcolor}
\usepackage{listings}

\lstset{basicstyle=\ttfamily\small,breaklines=true,frame=single,
        columns=fullflexible,keepspaces=true}

\title{\textbf{Experiment: AI-Assisted Integration of an Existing App into ShinySwarm}\\[0.5em]
\large Can an agent re-engineer the BioDT Honeybee use case? (working notes)}
\author{Iva Gunjaca\\
Master of Software Engineering, University of Amsterdam}
\date{\today}

\begin{document}
\maketitle

% --------------------------------------------------------------------------
\section*{Attribution / Disclaimer}
% --------------------------------------------------------------------------

The honeybee (Beekeeper pDT) application re-engineered in this experiment
originates with the \textbf{LTER LIFE project} (BioDT use case; repositories
\texttt{BioDT/biodt-shiny} and \texttt{BioDT/uc-pollinators}). It is used here
\textbf{with the permission of the LTER LIFE project} for the purpose of this
Master's thesis. All credit for the original honeybee/BEEHAVE application, its
scientific model, and its UI design belongs to the LTER LIFE / BioDT authors.
The artifacts produced in this experiment are re-engineered adaptations for the
ShinySwarm architecture with a \emph{simulated} compute backend; they are not
the original work and do not reproduce the real BEEHAVE scientific compute.

% --------------------------------------------------------------------------
\section{What this experiment is (and is not)}
% --------------------------------------------------------------------------

The feasibility note (\texttt{biodt\_feasibility.tex}) argued \emph{on paper}
that the deferred BioDT Honeybee use case could be brought into ShinySwarm. This
note records a different thing: an \textbf{experiment} in which an AI coding
agent was asked to \emph{actually perform} that integration inside the project
repository, using the existing ShinySwarm benchmark apps as the pattern. The
research interest is not the honeybee app per se but the \textbf{process}: how
far an agent gets at integrating a real, third-party scientific application into
an existing architecture, what it produces, and where it needs a human.

\textbf{Scope honesty.} The environment cannot run the real BEEHAVE NetLogo
container (\texttt{ghcr.io/biodt/beehave}), so the compute backend was
\emph{simulated} -- a stand-in service that honours the same input/output
contract. This is a deliberate, stated limitation: the experiment tests the
\emph{integration plumbing and UI re-use}, not the scientific compute, which the
feasibility note already showed is a solved, containerised problem.

% --------------------------------------------------------------------------
\section{Method}
% --------------------------------------------------------------------------

\begin{enumerate}[nosep]
  \item \textbf{Prompt.} A single natural-language request: read the thesis docs
    and assess/perform integrating the ``BioDT beehive app'' into ShinySwarm,
    simulating backend values. Architecture target: \emph{both} REST and Kafka.
  \item \textbf{Grounding.} The agent read the existing repo: the REST and Kafka
    architectures, the eight benchmark apps (notably \texttt{map/} and
    \texttt{monte\_carlo/}), the Spring \texttt{CollaborationController} state
    relay, and the \texttt{statesnap} package.
  \item \textbf{Source re-use (key step).} Rather than inventing a lookalike UI,
    the agent cloned the \emph{real} upstream app \texttt{BioDT/biodt-shiny} and
    read the actual honeybee module (\texttt{app/view/honeybee/*},
    \texttt{app/logic/honeybee/*}) to extract the true input parameters and the
    output result-table column layout.
  \item \textbf{Build.} The agent wrote the integration artifacts (below) for
    both architectures and registered the new app in the Spring backends,
    \texttt{docker-compose}, and \texttt{nginx}.
  \item \textbf{Verification.} R was installed in the environment; all four R
    artifacts were parse-checked, and the simulation model was executed to
    confirm it produces a BioDT-shaped, non-deterministic result table.
\end{enumerate}

% --------------------------------------------------------------------------
\section{What the agent produced (artifacts)}
% --------------------------------------------------------------------------

\begin{table}[h]
\centering\small
\begin{tabular}{@{}lp{8.6cm}@{}}
\toprule
\textbf{Path} & \textbf{Role} \\
\midrule
\texttt{.../RESTAPI/shiny\_services/honeybee/shiny\_front\_beehave.r} &
  BioDT honeybee UI (map, params, lookup, plot) flattened to plain Shiny,
  wired to the REST relay + \texttt{statesnap}. \\
\texttt{.../RESTAPI/shiny\_services/honeybee/shiny\_back\_beehave.r} &
  Plumber service; simulated BEEHAVE honouring BioDT's CSV/result contract. \\
\texttt{.../Kafka/shiny\_services/honeybee/shiny\_front\_beehave.r} &
  Same UI, event-driven (Producer/Consumer) collaboration. \\
\texttt{.../Kafka/shiny\_services/honeybee/shiny\_back\_beehave.r} &
  Kafka consumer/producer worker running the same model. \\
\texttt{CollaborationController.java} (REST) &
  One line: \texttt{"Honeybee" -> http://shiny-back-beehave:8000/state}. \\
\texttt{DataInitializer.java} (both) &
  App \#9 seeded (IDs 1--8 preserved so k6 benchmarks stay valid). \\
\texttt{docker-compose.yml}, \texttt{nginx.conf} (both) &
  Front+back service entries, proxy port, route. \\
\bottomrule
\end{tabular}
\caption{Artifacts produced by the agent in one session.}
\end{table}

\paragraph{Input contract reproduced verbatim.} The front sends BioDT's actual
parameter names from \texttt{beekeeper\_param.R}:
\texttt{N\_INITIAL\_BEES}, \texttt{N\_INITIAL\_MITES\_HEALTHY},
\texttt{N\_INITIAL\_MITES\_INFECTED}, \texttt{HoneyHarvesting},
\texttt{VarroaTreatment}, \texttt{DroneBroodRemoval}, plus the
year/\texttt{DaysLimit} simulation window and the editable habitat lookup table.

\paragraph{Output contract reproduced verbatim.} The backend emits the exact
columns BioDT's \texttt{honeybee\_beekeeper\_plot.R} reads:
\texttt{date}, \texttt{weather},
\texttt{TotalIHbees + TotalForagers}, and
\texttt{(honeyEnergyStore / ( ENERGY\_HONEY\_per\_g * 1000 ))}.

% --------------------------------------------------------------------------
\section{What ShinySwarm adds that BioDT lacked}
% --------------------------------------------------------------------------

\begin{itemize}[nosep]
  \item \textbf{Real-time collaboration.} Apiary coordinates, edited lookup-table
    cells, and run results are shared across users through the relay (REST poll
    at 500\,ms) or Kafka topics (\texttt{input}/\texttt{output}, keyed by
    session). Upstream BioDT is one-container-per-user ShinyProxy with no
    cross-user sharing.
  \item \textbf{Role/permission gating} reused from the benchmark apps
    (VIEWER vs EDITOR/OWNER), driving both UI disabling and producer creation.
  \item \textbf{Reproducible checkpoints} via \texttt{statesnap}: because the
    BEEHAVE model is \emph{non-deterministic}, input-only Shiny bookmarking
    cannot reproduce a colleague's exact run. The agent captured the
    \emph{computed result} (\texttt{state\_rds}) alongside inputs, so restore
    reproduces the exact trajectory -- the cross-user reproducibility claim the
    proposal made, now on a real scientific app.
\end{itemize}

% --------------------------------------------------------------------------
\section{Verification results}
% --------------------------------------------------------------------------

\begin{itemize}[nosep]
  \item \textbf{Syntax.} All four R artifacts parse cleanly (\texttt{R 4.3.3},
    \texttt{parse()}); the \texttt{statesnap} sources parse cleanly.
  \item \textbf{Execution.} The simulation function runs and returns a
    $366\times4$ table for a 365-day run, with BioDT's column names.
  \item \textbf{Non-determinism confirmed.} Two runs with different seeds gave
    peak colony sizes of $18{,}995$ vs $19{,}323$ bees -- this is precisely the
    property that makes the \texttt{statesnap} reproducibility argument non-trivial.
  \item \textbf{Not yet run end-to-end.} The full Docker stack (Postgres, Redis,
    Spring, Kafka, Angular) was \emph{not} brought up in this environment; the
    Java/compose/nginx wiring is verified by inspection and convention-matching,
    not by a live integration test. This is the obvious next validation step.
\end{itemize}

% --------------------------------------------------------------------------
\section{Where the agent needed a human / open issues}
% --------------------------------------------------------------------------

\begin{itemize}[nosep]
  \item \textbf{Decisions, not code.} The human chose scope (full PoC, both
    architectures) and the simulate-the-backend constraint. The agent executed.
  \item \textbf{Rhino $\to$ plain Shiny flattening} was done by hand-translation:
    the agent collapsed BioDT's nested \texttt{box}/\texttt{moduleServer} module
    tree into a single flat app. This is faithful to the benchmark style but is
    \emph{not} the ``integrate in place'' option, and it sidesteps the
    \texttt{statesnap} module-namespace gap flagged in the feasibility note (the
    flat app has no nested namespaces, so capture works as-is).
  \item \textbf{Map fidelity.} BioDT's input map is a land-use raster with
    in/out-of-boundary checks (\texttt{terra::extract}); the PoC uses a plain
    basemap + draw tool. The collaboration/contract is faithful; the
    geospatial validation is simplified.
  \item \textbf{Licensing} (\texttt{biodt-shiny} = ``Other'') must be confirmed
    before redistributing adapted UI code -- unchanged from the feasibility note.
\end{itemize}

% --------------------------------------------------------------------------
\section{Reading for the thesis (so what?)}
% --------------------------------------------------------------------------

As an integration-process result: an agent, given only the repo and a one-line
brief, (i) located and read the \emph{real} upstream app, (ii) recovered its
exact I/O contract, (iii) re-used the existing ShinySwarm collaboration pattern
rather than inventing one, and (iv) produced parseable, contract-faithful
artifacts for \emph{both} architectures plus the backend wiring, in a single
session. The residual work is the same scoped, known work the feasibility note
predicted (live end-to-end run; real BEEHAVE container; raster map fidelity;
in-place Rhino integration with namespace-aware \texttt{statesnap}). The
experiment therefore supports a narrow but useful claim for the discussion: the
ShinySwarm architecture is \emph{regular enough} that adding a ninth, real-world
app is largely mechanical -- mechanical enough that an agent can do the bulk of
it -- which is itself evidence for the architecture's extensibility.

% --------------------------------------------------------------------------
\section{Provenance}
% --------------------------------------------------------------------------

\begin{itemize}[nosep]
  \item Upstream read on 2026-06-12: \url{https://github.com/BioDT/biodt-shiny}
    (\texttt{app/view/honeybee/*}, \texttt{app/logic/honeybee/*}).
  \item Patterns re-used from this repo: \texttt{shiny\_services/map},
    \texttt{shiny\_services/monte\_carlo}, \texttt{statesnap/},
    \texttt{shinyswarm/.../session/CollaborationController.java}.
  \item All artifacts committed under the two architecture directories; see the
    table in \S4.
\end{itemize}

\end{document}
